\documentclass[11pt,a4paper]{article}

\usepackage[margin=2.5cm]{geometry}
\usepackage{setspace}
\usepackage{authblk}
\usepackage{amsmath,amssymb}
\usepackage{enumitem}
\usepackage{url}
\usepackage[round]{natbib}
\usepackage{hyperref}

\setstretch{1.15}
\parskip=6pt
\parindent=0pt
\newcommand{\PG}{P(\mathrm{green})}

\title{\textbf{Rebuilding-time based reference points for non-retained bycatch:\\
a forward-looking risk-equivalent example for North Atlantic shortfin mako}\\[0.5em]
\large Journal manuscript, written after SCRS/2026/XXX is withdrawn}

\author[1]{L. T. Kell}
\affil[1]{Centre for Environmental Policy, Imperial College London, United Kingdom}

\date{}

\begin{document}
\maketitle

\begin{abstract}
\noindent\textit{Draft abstract.}
Uncertainty methods for stock \emph{status} (delta method, bootstrap, MCMC)
are well compared \citep{Magnusson2013}, and the broader forecasting agenda
was set by \citet{Patterson2001}. Modern RFMO advice, however, is often a
\emph{projection}: Kobe strategy matrices, recovery probabilities and
rebuilding times under constant catch or harvest control rules. Those
quantities inherit process error, age-structured lags and structural
disagreement among models. Using North Atlantic shortfin mako as the primary
case---with supporting diagnostics and backtests on other stocks---we (i)~show
that multivariate log-normal (MVLN) Kobe probabilities are closed-form and need
no Monte Carlo; (ii)~compare MVLN, MCMC and age-structured process-error
projections under consistent recruitment assumptions; (iii)~contrast
age-structured SS3 and a JABBA3 surplus-production tier
(\texttt{ss3diags2::ss2jabba3}) on the risk-equivalent catch of
\citet{Fischer2023} and the rebuilding year of \citet{KellAcceptedFaF}, in
light of cryptic spawners and exploitable--SSB lags \citep{Winker2020SMA};
(iv)~propose diagnostics that
gate ensemble membership \citep{Kell2026Diags,Carvalho2021}; (v)~outline
objective selection/exclusion/weighting across frameworks
\citep{WangEnsembleR5}; and (vi)~argue that rebuilding targets need not be
MSY-based for sharks \citep{KellAcceptedFaF}. The projection record is
submitted first as an ICCAT SCRS paper and, after withdrawal, attached as
the supplementary material of this manuscript. This manuscript is written
after that withdrawal and does not repeat the grid or the tables.
\end{abstract}

\textbf{Keywords:} fisheries uncertainty; MVLN; MCMC; process error; Kobe;
ensemble; rebuilding; shortfin mako; diagnostics; backtesting

\section{Relationship to the SCRS document}

The order of writing is fixed.

\begin{enumerate}[leftmargin=*]
  \item \texttt{scrs\_2026\_proj\_supplement.tex} is the ICCAT paper. Its
    introduction and discussion are the top and tail around the projection
    record: four SMA-N projections, one catch grid, the trajectories, the
    Kobe matrices, the catches at $\PG = 0.6$ and $0.7$, the rebuilding
    years, and the closed-form check against the published matrix.
  \item That paper is submitted to the SCRS, then withdrawn.
  \item The withdrawn file is attached, unchanged, as the supplementary
    material of this manuscript. The design, the four specifications, the
    results and Appendix~A are not rewritten for the journal.
  \item This manuscript is written after the withdrawal. It keeps the
    forward-looking risk-equivalent argument, the status-versus-projection
    framing, diagnostics and ensemble weighting, and the case for rebuilding
    targets that need not be MSY-based for a non-retained bycatch
    \citep{SharkWG2026,KellAcceptedFaF}. It does not repeat the projection
    machinery.
\end{enumerate}

The earlier prospectus and its three split stubs are in \texttt{archive/}.
Outlines for papers not yet written are in \texttt{draft\_outlines.tex}.

\section{Introduction (planned)}

\begin{enumerate}[leftmargin=*]
  \item \textbf{Status vs projection.} \citet{Magnusson2013} on delta /
    bootstrap / MCMC for status; \citet{Patterson2001} on assessment and
    forecasting; gap for Kobe-style projection advice.
  \item \textbf{Why sharks expose the gap.} \citet{Winker2020SMA}: cryptic
    mature biomass, lag exploitable$\leftrightarrow$SSB, SPM misinterpretation;
    2026 SMA assessment status mix vs SS3-only projections
    \citep{SharkWG2026}.
  \item \textbf{Questions.} When is closed-form MVLN enough? When do MCMC and
    PE change advice? How should ensembles be built? What rebuild metric if
    not BMSY?
\end{enumerate}

\section{Materials and methods (planned)}

\subsection{Case studies}
\begin{itemize}[leftmargin=*]
  \item \textbf{Primary:} North Atlantic shortfin mako 2026 (SS3 MVLN and
    MCMC via \texttt{ss3diags2}; JABBA3 via \texttt{ss2jabba3}; FLasher PE)
    --- SCRS numerical core.
  \item \textbf{Advice-linked:} Scenario~1 / HW2final Kobe path
    (\texttt{sma-ll-2026}).
  \item \textbf{Diagnostics:} three-strand examples
    \citep{Kell2026Diags}; journal strand in \texttt{flrpapers/3strands}
    (pandalus, herring, sma).
  \item \textbf{Backtest:} historical peel projections vs later outcomes
    (\texttt{flrpapers/backtest}) as forecast-skill checks where claimed.
\end{itemize}

\subsection{Uncertainty treatments}
MVLN (analytical $\PG$); MCMC; circular-block recruitment PE in age-structured
forward projection; surplus-production posteriors. Explicit design matrix
attributing method vs structure vs approximation error (as in SCRS).

\subsection{Ensembles and diagnostics}
Gate runs by diagnostics before weighting \citep{Kell2026Diags,Carvalho2021};
weighting recipes from \citet{WangEnsembleR5}. Do not average SS3 and JABBA3
projections for SMA advice without that gate (WG chose SS3-only projections).

\subsection{Rebuild metrics}
Time-to-target and risk-equivalent catches \citep{KellAcceptedFaF}; sensitivity
to target definition (Kobe green vs SPR / depletion).

\section{Results (map from SCRS + extensions)}

\begin{itemize}[leftmargin=*]
  \item Analytical MVLN vs simulated KSM (SCRS Table~Liniers).
  \item Status and projection panels: SS3 MVLN, MCMC, FLasher PE, JABBA3.
  \item Risk-equivalent catch at $\PG = 0.6$ and $0.7$, and rebuilding year,
    as the open-loop form of \citet{Fischer2023}.
  \item PE vs deterministic: smooth SSB, fatter risk near KSM cliff.
  \item Rebuild-year curves; catch at 60/70\% $\PG$ by 2070.
  \item \textit{Journal adds:} backtest skill tables; SPR-target sensitivity;
    multi-stock diagnostic appendix; ensemble weight scenarios.
\end{itemize}

\section{Discussion (planned themes)}

\begin{enumerate}[leftmargin=*]
  \item Update Patterson/Magnusson for the projection era.
  \item Process error matters most for prediction, not only reconstruction.
  \item Age structure as both biology (Winker) and filter of PE.
  \item Ensembles need procedures, not informal hedging.
  \item Feedback management vs single open-loop projection to 2070.
  \item Submit the SCRS paper, withdraw it, and attach that file as the
    supplement. Write this manuscript afterwards. Do not expand the SCRS
    paper into the journal article.
\end{enumerate}

\section{Conclusions (planned)}

Operational checklist for RFMOs: analytical MVLN; consistent-recruitment
MCMC/PE checks; diagnostic gates; documented ensemble rules; rebuild metrics
matched to life history (not default BMSY for poorly aged sharks).

\section*{Acknowledgements}
\textit{To be completed.}

\bibliographystyle{plainnat}
\bibliography{refs}

\end{document}
